\section{Task and Dataset}
\label{sec:task}

\subsection{Task Definition}

ElCardioCC 2026 is a multi-label ICD-10 coding shared task for Greek cardiology
discharge summaries, organised as part of BioASQ at CLEF
2026~\cite{BioASQ2026,BioASQ2026ElCardioCC}, with 115 target codes.

\subsubsection{Evaluation metric.}
The primary metric is a group-level F1 score that reflects clinical equivalence among
related codes. Each gold annotation is a \emph{group}, a set of interchangeable
ICD-10 codes. Scoring is defined as follows:
\begin{itemize}
  \item TP: a gold group is matched if at least one of its member codes is predicted.
  \item FN: gold groups with no predicted member.
  \item FP: predicted codes that do not belong to any gold group.
\end{itemize}
Multiple predictions within the same group do not increase TP. Precision, recall, and
F1 are then:
\begin{equation}
P = \frac{\text{TP}}{\text{TP}+\text{FP}}, \quad
R = \frac{\text{TP}}{\text{TP}+\text{FN}}, \quad
F1 = \frac{2PR}{P+R}
\end{equation}
This differs from standard multi-label micro-F1: predicting any one equivalent code
fully satisfies the corresponding group, whereas predicting codes outside any gold group
is penalised as false positives.

\noindent\textbf{Example.}
Given gold groups $\{\{$R55$\}, \{$I10, I11$\}, \{$I44$\}, \{$Z95$\}\}$:
\begin{itemize}
  \item Predicting \{R55, I10, I11, I44, Z95\}: TP=4, FP=0, F1=1.0; both I10 and I11 satisfy
    the same group without penalty, since neither falls outside a gold group.
  \item Predicting \{R55, I10, I44, Z95, I20\}: TP=4, FP=1, F1=0.889; I20 does not belong
    to any gold group and is counted as a false positive.
\end{itemize}

\subsection{Dataset}

The dataset consists of 2,500 anonymised discharge summaries. We split these into
training (2,000), validation (250), and test (250) sets; a separate blind test set is
held by the organisers. Documents are moderately long: mean 280 tokens ($\approx$2,000 characters),
maximum 1,218 tokens. Only 2.4\% of documents exceed the 512-token BERT limit.

The 115 codes are highly skewed. The five most frequent (Z95, I25, I21, I48, E11)
account for $\sim$41\% of all annotations; 30 codes appear fewer than 10 times and 4
codes appear only once (Table~\ref{tab:label_stats}). This long-tailed distribution
requires explicit strategies for rare-label coverage.

\begin{table}[ht]
\centering
\caption{Label frequency distribution across 115 ICD-10 codes.}
\label{tab:label_stats}
\begin{tabular}{lrr}
\toprule
Frequency band & \# codes & Approx.\ \% of instances \\
\midrule
$\geq 500$  & 5  & 41\% \\
100--499    & 22 & 36\% \\
10--99      & 58 & 21\% \\
$< 10$      & 30 & $<$2\% \\
\bottomrule
\end{tabular}
\end{table}

\subsection{Preprocessing and Data Augmentation}

Text cleaning is shared across all splits: whitespace collapse and removal of non-essential
special characters (preserving hyphens, commas, slashes, parentheses). Greek BERT
additionally receives lowercasing and accent stripping; XLM-RoBERTa receives
whitespace-only normalisation to preserve abbreviation signals (``PCI'', ``TAVI'').
ICD-10 annotations are encoded as 115-dimensional binary vectors.

Stratified splitting uses a two-stage Multilabel Stratified Shuffle Split~\cite{sechidis2011stratification}
to produce an 80/10/10 partition, ensuring every validation and test label also appears
in training. Of the 2,500 samples, 1,453 include mention-level annotations enabling
NER+EL training.

Synonym augmentation (XLM-RoBERTa base only) draws directly on the dictionary
baseline's term-to-code table (Section~\ref{subsec:dict}): for each gold code on a
document with at least two known surface forms, one occurrence of a matched term is
replaced, with probability $p=0.35$, by another randomly chosen term mapped to the
same code. Frequency-adaptive multipliers generate additional augmented copies per
document ($\times$4 for codes with 30--99 training instances, $\times$3 for
100--199, $\times$2 for 200--299; codes with fewer than 30 or at least 300 instances
are not augmented). For example, for code I21 (myocardial infarction), a discharge
summary sentence translating to ``\emph{Diagnosis: acute myocardial infarction
(STEMI) of the inferior wall}'' (original in Greek) may have the matched term
``STEMI'' replaced by an alternative I21 surface form such as ``NSTEMI'' drawn from
the same term list, yielding ``\emph{Diagnosis: acute myocardial infarction
(NSTEMI) of the inferior wall}''; because substitution is a case-insensitive regex
replacement rather than a grammar-aware rewrite, it can occasionally produce minor
casing or spacing artefacts around the substituted span, which we treat as an
acceptable trade-off given the resulting recall gains on underrepresented codes.
